% TOP_PROBLEM: {"id": 2634, "title": "Kourovka Notebook Problem 21.125", "category_id": 17, "category": {"id": 17, "name": "group_theory", "display_name": "Group Theory", "description": "Problems about groups, group actions, representations, and related algebraic structures.", "slug": "group-theory", "order_index": 17, "created_at": "2026-07-31T15:26:25.671Z"}, "status": "open", "problem_number": "KOU-21.125", "set_id": 10}
% FIRST_POSTED: 2026-10-01
% ENTRY_KIND: statement_only
% UnsolvedMath ID 2634; source catalogue code KOU-21.125
% !TeX program = lualatex
% Definition and sources only; this file is not an LLM solution attempt.
\documentclass[11pt,a4paper]{article}
\usepackage[margin=25mm]{geometry}
\usepackage{fontspec}
\setmainfont{lmroman10-regular.otf}[BoldFont=lmroman10-bold.otf,
 ItalicFont=lmroman10-italic.otf,BoldItalicFont=lmroman10-bolditalic.otf]
\usepackage{amsmath,amssymb,mathtools}
\usepackage{unicode-math}
\setmathfont{latinmodern-math.otf}
\AtBeginDocument{\let\setminus\smallsetminus}
\usepackage{xurl,hyperref}
\hypersetup{colorlinks=true,urlcolor=blue}
\setcounter{secnumdepth}{0}
\setlength{\parindent}{0pt}
\setlength{\parskip}{5pt}
\setlength{\emergencystretch}{3em}
\begin{document}
\section*{Kourovka Notebook Problem 21.125}\label{2634}
{\small UnsolvedMath ID 2634 · KOU-21.125 · Group Theory · Kourovka Notebook - New Problems, Issue 21}
\subsection{Definitions and mathematical statement}
Let $F_m$ be the free group of rank $m\ge1$ with basis $x_1,\ldots,x_m$, and let $|g|$ be reduced word length in that basis. Suppose $\varphi\in\operatorname{Aut}(F_m)$ satisfies
\[
 c n^{m-1}\le\max_i|\varphi^n(x_i)|\le C n^{m-1}
 \quad(n\ge n_0)
\]
for some $c,C>0$ and $n_0\ge1$. This specifies polynomial growth of degree $m-1$ as in the source. Define the free-by-cyclic group
\[
 G_\varphi=\langle x_1,\ldots,x_m,t\mid
 t^{-1}x_it=\varphi(x_i)\ (1\le i\le m)\rangle.
\]

For completeness, a cube complex is formed by gluing faces of unit Euclidean cubes by isometries. A Salvetti complex is the union of coordinate tori indexed by the finite cliques of a graph, with one common vertex and one generator loop per graph vertex. A cube complex is special if it admits a cubical local isometry to such a Salvetti complex. On vertex links this means an injective map onto a full subcomplex; full means that every simplex whose vertices lie in the image also lies in the image. A group is virtually special if some finite-index subgroup is isomorphic to the fundamental group of a special cube complex.

\textbf{(a)} Is $G_\varphi$ virtually special for every automorphism satisfying the displayed growth hypothesis?

\textbf{(b)} In particular, is every Hydra group
\[
 G_m=\langle a_1,\ldots,a_m,t\mid
 t^{-1}a_1t=a_1,\quad t^{-1}a_it=a_ia_{i-1}\ (2\le i\le m)\rangle
\]
virtually special? The target is existence of some finite-index subgroup with the stated cubical realization. No bound on its index, dimension, or number of cubes is required in either part.
\subsection{Short English statement}
Are free-by-cyclic groups with maximal-degree polynomial monodromy virtually special, including all the explicitly presented Hydra groups?
\subsection{Sources}
\begin{itemize}
\item[S1] ULAM.AI and the UnsolvedMath Contributors, supplied problems.json, record 2634 (KOU-21.125), Kourovka Notebook - New Problems, Issue 21. Catalogue identity and source transcription; CC BY 4.0.
\url{https://huggingface.co/datasets/ulamai/UnsolvedMath}.
\item[S2] The Kourovka Notebook, 21st edition, new problems, Problem 21.125. Expanded formulation of the supplied catalogue statement.
\url{https://arxiv.org/pdf/1401.0300v47}.
\item[S3] N. Brady and I. Soroko, Dehn functions of subgroups of right-angled Artin groups, Theorem 4.12; local-isometry criterion for special cube complexes.
\url{https://sites.math.unt.edu/~soroko/Dehn-functions-RAAGs.pdf}.
\end{itemize}
\end{document}
